% Introduction (input by main.tex after \section{Introduction})

\paragraph{Three notions of causality.}
When a computer scientist says that one part of a system \emph{causes} another, they may mean at least three distinct things. In Pearl's theory of causal inference \cite{pearl2009causality}, $X$ has a causal effect on $Y$ if \emph{intervening} on $X$---forcing it to a value, regardless of how it would otherwise have been produced---changes the distribution of $Y$. In probabilistic modelling, a causal structure is a directed acyclic graph along which a joint distribution \emph{factorises}, and the content of the structure is the family of conditional independences it entails \cite{pearl2009causality,fritz2023dseparation}. In quantum foundations and in the theory of higher-order quantum processes \cite{kissinger2019categorical,oreshkov2012noncausal}, the causal structure of a process is given by its \emph{signalling} constraints, which inputs can influence which outputs, and in concurrency theory a causal structure is an explicit partial order on events \cite{lamport1978time,pratt1986pomsets}.

Three further usages should be named so as to be set aside. In categorical quantum mechanics, a process is called \emph{causal} when it preserves the discarding map \cite{PQP_Coecke_Kissinger_2017}; this is the causality axiom of Chiribella, D'Ariano and Perinotti \cite{chiribella2010purification}, glossed as ``no signalling from the future'', and Coecke \cite{coecke2014terminality} shows that it implies the no-signalling principle for processes composed along a causal structure; it is the base axiom on which signalling structure is built, not a notion of causality between parts. Fritz \cite{fritz2020synthetic} calls a Markov category \emph{causal} when the output of a morphism cannot depend on a choice made downstream of it; this is again a property of the ambient category. In artificial intelligence and in verification, \emph{actual causality} \cite{halpern2016actual} asks which particular event caused which particular outcome in a single run; that question is not treated here. Nor are experimental tests of indefinite causal order or causal discovery from data. Identifiability and the do-calculus are not set aside: the proposal reaches them (\S\ref{sec:prop-graph}).

That the same word carries these meanings is not an accident, but neither is it a theorem. Each notion is \emph{extensional}: it characterises causal structure by what an observer may do to the system---intervene, condition, send signals. What such observers probe is \emph{intensional}: the dataflow of a program, the graph of a model, the wiring of a circuit. Programs carry this intensional structure explicitly in their syntax. The viewpoint of this thesis is therefore that causality is a property of \emph{program semantics}, and that the right way to study a notion of causality is to relate the intensional structure in the syntax to its extensional characterisation in the semantics by a correspondence theorem. A recurring theme, made precise in \S\ref{sec:synthesis}, is that the two never quite coincide: a dependence that is present in the syntax can be invisible to every observer, by what the causal-inference literature calls \emph{fine tuning}.

\paragraph{Where I stand.}
In my first year, with coauthors, I built a hierarchy of first-order probabilistic languages in which Pearl's interventions and counterfactual \emph{twinning} are program transformations, defined program equivalences contextually, and gave a graded denotational semantics that is fully abstract for concrete models and extends to a higher-order language (\S\ref{sec:intv-wherewestand}). This gives the interventional notion its extensional face: equivalence under all interventional contexts, which the graded semantics characterises exactly. What it does not characterise is the dataflow itself, which is strictly finer (\S\ref{sec:intv-wherewestand}). In the running language, a program reads an \emph{intervention label} $\ell$ through $\mathsf{applyIntv}\ \ell\ t$, which yields the label's value when it is set and $t$ otherwise; if a second variable $z$ does not depend on the result, every intervention at $\ell$ leaves the distribution of $z$ unchanged, and the contextual equivalence makes this an equation between programs. Causal independence is \emph{tested by contexts}, and, in programming-language terms, it is a \emph{noninterference} property \cite{goguen1982noninterference}. The thesis takes this work as its subject and builds outward from it: towards practice and expressivity, towards a better understanding of the constructions it rests on, towards a graph-centred account, and, open-endedly, towards the notions of causality used in concurrency and in higher-order quantum processes.

\paragraph{Structure of this document.}
Sections \ref{sec:foundations}--\ref{sec:order} review the literature. Section \ref{sec:foundations} covers the semantic foundations: algebraic effects, probabilistic programs, Markov categories and grading. Section \ref{sec:intv} covers causal models, their equivalences, cyclic models, identifiability, their diagrammatic presentations, and causal probabilistic programming, which is the state of the art the proposal extends. Section \ref{sec:order} positions the two neighbouring fields that the proposal reaches into, higher-order quantum processes and the semantics of concurrency, at the level of their main definitions and results; their detailed survey is a milestone of the plan. Section \ref{sec:synthesis} draws the notions together. Section \ref{sec:proposal} is the proposal, organised in four directions and followed by evaluation, risks and the plan of work.
